\section{Introducción}
\subsection{Conceptos básicos generales}
\begin{frame}{Definiciones}
\begin{textblock*}{\textwidth}(0.4cm,1.2cm)
\begin{description}[<+(1)->] [Espacio de Contexto de ejecución:] \justifying %\parskip 0.3cm
 \item[Tarea:] Unidad de trabajo que un procesador puede despachar, ejecutar, y detener a voluntad, bajo la forma de:
 \begin{itemize} \justifying %\parskip 0.3cm
  \item Instancia de ejecución de un programa (\textit{proceso}, o \textit{thread} para el Sistema Operativo).
  \item Un handler de interrupción, o excepción.
  \item Un servicio del S.O. (por ejemplo en Linux, cualquiera de las Syscall).
 \end{itemize}
% \end{description}
% \end{textblock*}
% \end{frame}

% \begin{frame}{Definiciones}
%  \begin{textblock*}{\textwidth}(0.4cm,1.2cm)
%   \begin{description}[<+(1)->] [Espacio de Contexto de ejecución:] \justifying \parskip 0.3cm 
   \item [Espacio de ejecución:] Es el conjunto de secciones de código, datos, y pila que componen la tarea. 
   \item [Contexto de ejecución:] Es el conjunto de valores de los registros internos del procesador en cada momento. 
   \item [Espacio de Contexto de ejecución:] Bloque de memoria en el que el S.O. almacenará el contexto completo de ejecución del procesador.
  \end{description}
 \end{textblock*}
\end{frame}


\begin{frame}{Context Switch}
 \begin{textblock*}{\textwidth}(0.4cm,1.2cm)
  \begin{block}<1->{Definición} \justifying
Salvar y restaurar el estado computacional o contexto de dos procesos o threads cuando se suspende la ejecución del primero (se salva su contexto) para pasar a ejecutar el segundo (se restaura su contexto). Puede incluir, o no, el resguardo y restauración del espacio de memoria (procesos o threads).
  \end{block}
%  \end{textblock*}
% \end{frame}

% \begin{frame}{Context Switch}
%  \begin{textblock*}{\textwidth}(0.4cm,1.2cm)
  \begin{block}<2->{Primeras conclusiones} \justifying
Las tareas en un computador no operan simultáneamente, sino serialmente pero conmutando de una a otra a gran velocidad. Esto significa que se producen cientos o miles de context switches por segundo.\\
Nuestros sentidos no captan la intermitencia de cada tarea, creándose una sensación de simultaneidad.\\
  \end{block}
 \end{textblock*}
\end{frame}

\begin{frame}{Entra en juego el Sistema Operativo}
 \begin{textblock*}{\textwidth}(0.4cm,1.2cm)
  \begin{itemize}[<+(1)->] \justifying \parskip 0.3cm
   \item En particular, el Sistema Operativo tiene un módulo de software llamado \textbf{\textit{scheduler}}, que trabaja con una lista de tareas a ejecutar.
   \item El \textbf{\textit{scheduler}} define un intervalo de tiempo llamado time frame, el que a su vez con ayuda de un temporizador es dividido en intervalos mas pequeños, que se convierten en la unidad de tiempo.
   \item El arte de manejo de prioridades consiste entonces en asignarle a cada tarea un porcentaje del time frame, medido en una cantidad de intervalos unitarios.
   \item Cada tarea tiene así unos milisegundos para progresar, expirados los cuales suspende una tarea y despacha para su ejecución la siguiente de la lista.
  \end{itemize}
 \end{textblock*}
\end{frame}

\begin{frame}{Integrando lo visto hasta aquí}
 \begin{textblock*}{.95\textwidth}(0.4cm,1.1cm)
  \centering \includegraphics[width=\textwidth]{GenericTaskSwitch1.pdf}\\
 \end{textblock*}
 \begin{textblock*}{\textwidth}(0.4cm,7.1cm)
  \justifying
  \begin{small}
   \only <2->{Durante un tiempo T1, el procesador ejecuta la tarea \textit{Task1} de la lista que maneja el \textbf{\textit{scheduler}}.\\}
   \only <3->{El tiempo T1, que se le asignó para ejecución es medida de la prioridad de esta tarea.}
  \end{small}
 \end{textblock*}
\end{frame}

\begin{frame}{Integrando lo visto hasta aquí}
 \begin{textblock*}{.95\textwidth}(0.4cm,1.1cm)
  \centering \includegraphics[width=\textwidth]{GenericTaskSwitch2.pdf}\\
 \end{textblock*}
 \begin{textblock*}{\textwidth}(0.4cm,7.1cm)
 \justifying
  \begin{small}
   \only <2->{Una vez expirado T1, el \textbf{\textit{scheduler}} toma el control del sistema y como la tarea \textit{Task1 }no finalizó, la suspende, asignándole estado \textbf{\color{red!90!black}SUSPENDING}.\\}
   \only <3->{Resguarda el contexto de \textit{Task1} en un área de \textbf{memoria del kernel} (Flujo 1).\\}
  \end{small}
 \end{textblock*}
 \begin{textblock*}{.5\textwidth}(5.4cm,5.0cm)
  \only <4->{
   \begin{tcolorbox}[enhanced jigsaw,rounded corners, drop fuzzy shadow=black!50!white]
   \textbf{Contexto: Valor de los registros del procesador al momento de suspender la tarea.}
  \end{tcolorbox}
  }
 \end{textblock*}
\end{frame}

\begin{frame}{Integrando lo visto hasta aquí}
 \begin{textblock*}{.95\textwidth}(0.4cm,1.1cm)
  \centering \includegraphics[width=\textwidth]{GenericTaskSwitch3.pdf}\\
 \end{textblock*}
 \begin{textblock*}{\textwidth}(0.4cm,7.1cm)
  \justifying
  \begin{small}
   \only <2>{A partir del inicio de la transferencia del contexto de \textit{Task1} a la \textbf{memoria del kernel} inicia un intervalo que denominaremos $t_{c}$ (tiempo de conmutación).\\}
   \only <3->{Seguidamente el \textbf{\textit{scheduler}} busca la siguiente tarea en la lista de tareas en ejecución (en el gráfico \textit{Task2}), identifica la dirección de memoria de kernel en donde está almacenado su contexto, y marca a la tarea \textbf{\color{yellow!90!red}RUNNING}.}
  \end{small}
 \end{textblock*}
\end{frame}

\begin{frame}{Integrando lo visto hasta aquí}
 \begin{textblock*}{.95\textwidth}(0.4cm,1.1cm)
  \centering \includegraphics[width=\textwidth]{GenericTaskSwitch4.pdf}\\
 \end{textblock*}
 \begin{textblock*}{\textwidth}(0.4cm,7.1cm)
  \justifying
  \begin{small}
   \only <2-3>{Se carga en el procesador el contexto de la tarea \textit{Task2} (Flujo 2).\\}
   \only <3-3>{En el siguiente ciclo de clock, se reasume la ejecución de \textit{Task2}, en la instrucción siguiente a la última ejecutada antes de ser suspendida.\\}
   \only <4-5>{Si \textit{Task2} fue recién insertada en la lista de ejecución, entonces el procesador ejecutará su primer instrucción.\\}
   \only <5-5>{Llegamos al final del intervalo que denominamos $t_{c}$.}
  \end{small}
 \end{textblock*}
\end{frame}

\begin{frame}{Integrando lo visto hasta aquí}
 \begin{textblock*}{.95\textwidth}(0.4cm,1.1cm)
  \centering \includegraphics[width=\textwidth]{GenericTaskSwitch5.pdf}\\
 \end{textblock*}
 \begin{textblock*}{\textwidth}(0.4cm,7.1cm)
 \justifying
  \only <2->{El procesador ejecuta \textit{Task2} durante el tiempo T2.\\}
  \only <3->{\underline{\textbf{Nota}}: El Sistema Operativo asigna prioridades diferentes a cada tarea $\therefore T2 != T1$.\\}
  \only <4->{Nuevamente el \textbf{\textit{scheduler}} tomará el control cuando expire el tiempo T2.}
 \end{textblock*}
\end{frame}

\begin{frame}{Integrando lo visto hasta aquí}
 \begin{textblock*}{.95\textwidth}(0.4cm,1.1cm)
  \centering \includegraphics[width=\textwidth]{GenericTaskSwitch6.pdf}\\
 \end{textblock*}
 \begin{textblock*}{\textwidth}(0.4cm,7.1cm)
  \justifying
  \begin{small}
   \only <2-3>{El Sistema Operativo pone \textit{Task2} \textbf{\color{red!90!black}SUSPENDING} ya que su tiempo de ejecución expiró.\\}
   \only <3-3>{En el flujo 3, se ve el resguardo del contexto. Por otra parte se busca el identificador de la tarea siguiente en la lista.\\}
   \only <4->{Una vez ubicado, se identifica el puntero a la dirección de memoria en donde comienza el contexto de la siguiente tarea.}
  \end{small}
 \end{textblock*}
\end{frame}

\begin{frame}{Integrando lo visto hasta aquí}
 \begin{textblock*}{.95\textwidth}(0.4cm,1.1cm)
  \centering \includegraphics[width=\textwidth]{GenericTaskSwitch7.pdf}\\
 \end{textblock*}
 \begin{textblock*}{\textwidth}(0.4cm,7.1cm)
  \justifying
  \begin{small}
   \only <2-3>{Identificada la tarea siguiente, se carga su contexto en el procesador.\\}
   \only <3-3>{Observar que el tiempo $t_{c}$ es siempre el mismo independientemente de las tareas.\\}
   \only <4-5>{Finalizado este proceso en el ciclo de clock siguiente el procesador estará ejecutando la tarea \textit{Task3}.\\}
   \only <5->{Esta continuará durante un tiempo T3.}  
  \end{small}
\end{textblock*}
\end{frame}

\begin{frame}{Integrando lo visto hasta aquí}
 \begin{textblock*}{.95\textwidth}(0.4cm,1.1cm)
  \centering \includegraphics[width=\textwidth]{GenericTaskSwitch8.pdf}\\
 \end{textblock*}
 \begin{textblock*}{\textwidth}(0.4cm,7.1cm)
  \justifying
  \only <2->{Una vez expirado el tiempo de ejecución de \textit{Task3}, se repite el procedimiento con la siguiente tarea en la lista del \textbf{\textit{scheduler}}En este caso \textit{Task4}.\\}
  \only <3->{El flujo 4 muestra el resguardo del contexto de \textit{Task3}.}
 \end{textblock*}
 \end{frame}

\begin{frame}{Integrando lo visto hasta aquí}
 \begin{textblock*}{.95\textwidth}(0.4cm,1.1cm)
  \centering \includegraphics[width=\textwidth]{GenericTaskSwitch9.pdf}\\
 \end{textblock*}
 \begin{textblock*}{\textwidth}(0.4cm,7.1cm) \justifying
  \only <2->{El Flujo 5 muestra como se carga en el procesador el contexto de \textit{Task4}, que ejecutará durante el tiempo T4.}
 \end{textblock*}
\end{frame}

\begin{frame}{Integrando lo visto hasta aquí}
 \begin{textblock*}{.95\textwidth}(0.4cm,1.1cm)
  \centering \includegraphics[width=\textwidth]{GenericTaskSwitch10.pdf}\\
 \end{textblock*}
 \begin{textblock*}{\textwidth}(0.4cm,7.1cm)
  \justifying
  \only <2->{Finalmente llega el turno de la última tarea en la lista.\\}
  \only <3->{Una vez completada esta el \textbf{\textit{scheduler}} volverá al inicio de la lista y reasumirá la tarea \textit{Task1}.}
 \end{textblock*}
\end{frame}
